\chapter{Theoretical Foundation}
\label{chap:theory}

\section{Financial Data and Preprocessing}

\subsection{Transaction, Balance and Cash-Flow Model}

A minimally valid transaction has an amount $A>0$, a transaction type, a
date, a description, a category and a wallet within the user's scope.
PERFIN distinguishes \texttt{income}, \texttt{expense}, \texttt{transfer},
\texttt{investment\_inflow}, \texttt{investment\_outflow} and
\texttt{investment\_pnl} because each type has a different effect on
balances and reports. For a wallet $w$, the balance change caused by an
income or expense transaction is

\begin{equation}
  \Delta B_w=
  \left\{
  \begin{array}{ll}
    +A, & \mbox{if income},\\
    -A, & \mbox{if expense}.
  \end{array}
  \right.
\end{equation}

Over a period $P$, total income $I_P$, total expense $E_P$ and net cash
flow $N_P$ are computed from transactions that have not been soft-deleted:

\begin{equation}
  I_P=\sum_{t\in P,\,type(t)=income}A_t,\qquad
  E_P=\sum_{t\in P,\,type(t)=expense}A_t,\qquad
  N_P=I_P-E_P.
\end{equation}

Transferring $A$ from a source wallet $s$ to a destination wallet $d$
produces $\Delta B_s=-A$ and $\Delta B_d=+A$, so that
$\Delta(B_s+B_d)=0$. This invariant ensures a wallet transfer is not
counted as income or expense. The prototype allows a negative wallet
balance to reflect cash or an overdrawn account; the business rule is a
positive amount, two distinct wallets belonging to the same user, not a
``sufficient balance'' condition. The debit, credit and history record
must lie within the same \texttt{BEGIN/COMMIT} transaction; if one step
fails, \texttt{ROLLBACK} restores the entire state.

\subsection{Data Integrity and Lifecycle}

PERFIN applies four layers of protection:

\begin{enumerate}
  \item \textbf{Domain constraints:} transaction amounts, limits and
  contribution amounts must be valid; dates must not exceed the allowed
  business domain.
  \item \textbf{Referential integrity:} transactions reference users,
  categories and wallets by foreign keys.
  \item \textbf{Business integrity:} no transfer between the same wallet;
  no transaction written with a missing required field; soft delete does
  not remove the ability to recover.
  \item \textbf{Atomicity:} every multi-table change or balance change
  together with its history is committed in full or not at all.
\end{enumerate}

Indexes are placed on frequently filtered columns such as user, date,
category, wallet and status. Money is stored as
\texttt{DECIMAL(15,2)}; when passed to JavaScript it must be converted
in a controlled manner to avoid string concatenation or unintended
rounding. The cache is invalidated after commit. A cache error after
commit must not be returned as a business error, because the client may
retry and cause a duplicate effect.

\subsection{Time-Axis Normalization and Observation Windows}

Analysis by day, week or month must preserve the correct calendar axis.
Let $x_j$ be the total expense of the $j$-th calendar period in a window
$W$. If that period has no transactions then $x_j=0$; the period must not
be removed from the series. Zero-fill prevents two distortions:
compressing the time gaps inflates the slope, and taking only days with
spending inflates the burn rate.

The observation window must always accompany the result. PERFIN currently
uses 14 days for runway, 30 days for anomaly, 90 days for recurring items,
6 months for trend and 12 weeks for correlation by default. These are the
prototype's starting parameters, not thresholds proven optimal across all
users.

\subsection{Durable Data and Transient State}

PostgreSQL stores auditable business data. Redis stores short-lived state:
transactions awaiting confirmation, the field being clarified, the list of
ambiguous choices, caches and rate-limit counters. TTL expires
conversation state through the \texttt{EXPIRE} mechanism. The in-process
memory fallback serves development only; its data is lost on restart and
is not a production configuration.

\section{Input Normalization and Classification}

\subsection{Normalizing Amount, Date and Transaction Structure}

Natural input is converted into a typed draft before validation. The
suffixes \texttt{k}, ``nghìn'', ``ngàn'' multiply the value by $10^3$;
\texttt{tr}, ``triệu'' multiply by $10^6$. Relative dates are resolved
against the local date of the request. A sentence with several items must
produce an array of drafts, in which each element has its own
\texttt{amount}, \texttt{type}, \texttt{description},
\texttt{transaction\_date}, \texttt{category} and \texttt{wallet}. The
parser or LLM only creates the draft; the service then checks for a finite
number, the date domain, the category, the wallet and ownership before
creating the preview.

\subsection{Text Similarity and Safe Category Selection}

Strings are stripped of diacritics, lowercased, cleared of non-alphanumeric
characters and collapsed on whitespace. For two normalized strings $a,b$,
the edit score uses the Levenshtein distance $D_{lev}$~\cite{levenshtein1966}:

\begin{equation}
  s_{edit}=1-\frac{D_{lev}(a,b)}{\max(|a|,|b|)}.
\end{equation}

Let $T_a,T_b$ be two token sets; the Dice coefficient~\cite{dice1945} is

\begin{equation}
  s_{dice}=\frac{2|T_a\cap T_b|}{|T_a|+|T_b|}.
\end{equation}

If the smaller token set has at least two elements and is fully contained
in the larger set, the containment score $s_{contain}$ lies in the range
$[0{,}82;0{,}94]$ according to the length ratio. The overall score in the
source code is

\begin{equation}
  s=\max(s_{edit},\;0{,}92s_{dice},\;s_{contain}).
\end{equation}

The selection order is exact category name, exact alias, then fuzzy. Long
input uses the threshold $\tau=0{,}82$, input of no more than four
characters uses $\tau=0{,}90$. The best candidate must also beat the
second candidate by at least a margin $m=0{,}08$. If $s_1<\tau$ or
$s_1-s_2<m$, the system chooses ``Other'' or requests clarification rather
than assigning an ambiguous result on its own. This formula serves FR-04
and does not replace the user's confirmation.

\subsection{Learning from Feedback}

When the user corrects a category, the system stores the original
result--correct result pair together with the context. Exact/fuzzy
corrections are considered before aliases and the LLM on the next round; a
conflicting correction lowers confidence or requests confirmation. This is
retrieval combined with adaptive rules, \textbf{not model fine-tuning}.
The effect of feedback must be evaluated by accuracy or macro-F1 before and
after on the same replay set.

\subsection{OCR, Speech-to-Text and Provider Adapters}

OCR comprises image preprocessing, text-region detection, recognition and
post-processing. STT converts audio into a transcript. In PERFIN, the
provider returns only raw text or a transcript along with the provider name
and error status; the result then continues through the
draft--validation--preview pipeline as text.

A receipt image may produce one aggregate transaction or several line-item
drafts, but the current version \textbf{does not yet have an algorithm to
automatically reconcile} the line-item total with the invoice total, tax
and discount. The user must check the preview; a two-image smoke test is
not a measurement of OCR accuracy. The transcript must also be shown before
confirmation so that a recognition error does not silently become an
amount. The adapter allows a local provider to be replaced by a cloud
provider without modifying the transaction service. When a provider fails,
the system returns an error or allows manual entry; it does not generate
sample data.

\section{Analysis and Planning Algorithms}

\begin{table}[htbp]
  \centering
  \caption{Deterministic algorithms and the features that use them}
  \label{tab:algorithms}
  \small
  \begin{tabularx}{\textwidth}{|p{3.0cm}|X|X|}
    \hline
    \textbf{Algorithm} & \textbf{Input and output} &
    \textbf{Feature that uses it} \\ \hline
    Linear regression & Monthly expense series with zero-fill; slope, $R^2$,
    forecast & FR-07: trend \\ \hline
    Z-score and IQR & Daily or per-item expense series; flags, scores,
    method & FR-07: anomaly \\ \hline
    Cashflow runway & Total balance and 14 calendar days; burn rate, day
    count, depletion date & FR-07/FR-11: cash-flow alert \\ \hline
    Recurring items & Description, amount, date; clusters, cadence, monthly
    estimate & FR-07/FR-11: subscription scan \\ \hline
    Pearson & Two weekly series with zero-fill; positive correlation &
    FR-07: category relationships \\ \hline
    Budget recommender & Monthly expense, income, limits; recommendations
    and overrun forecast & FR-09 \\ \hline
    Goal planner & Goal, contribution or repayment, deadline, interest rate;
    time, shortfall, what-if & FR-10 \\ \hline
  \end{tabularx}
\end{table}

\subsection{Linear Regression and Trend Forecasting}

Let $y_i$ be the total expense of period $i$, $x_i=i$ with
$i=0,\ldots,n-1$, and $\bar{x}$ and $\bar{y}$ the corresponding means. The
least-squares (OLS) line $\hat y_i=a+bx_i$~\cite{montgomery2012} has

\begin{equation}
  b=\frac{\sum_{i=0}^{n-1}(x_i-\bar{x})(y_i-\bar{y})}
          {\sum_{i=0}^{n-1}(x_i-\bar{x})^2},\qquad
  a=\bar{y}-b\bar{x}.
\end{equation}

The goodness of fit and the next-period forecast are

\begin{equation}
  R^2=1-\frac{\sum_i(y_i-\hat y_i)^2}
                 {\sum_i(y_i-\bar y)^2},\qquad
  \hat y_n=\max(0,a+bn).
\end{equation}

When a constant series makes the denominator of $R^2$ equal to 0, the
source code returns $R^2=0$ and does not report a trend. The service adds
a rising trend to the facts only when there are at least three months,
$b>0$, $R^2\geq0{,}5$ and the average percentage change between pairs with
a positive prior sample reaches at least 10\%. Zero-fill must be performed
before the regression. The result describes a short-term linear trend; it
does not model seasonality or future events.

\subsection{Anomaly Detection with Z-score and IQR}

For $n\geq4$ expense values $x_i$, the mean, sample standard deviation and
z-score are

\begin{equation}
  \mu=\frac{1}{n}\sum_{i=1}^{n}x_i,\qquad
  s=\sqrt{\frac{\sum_{i=1}^{n}(x_i-\mu)^2}{n-1}},\qquad
  z_i=\frac{x_i-\mu}{s}.
\end{equation}

Quantiles are linearly interpolated over the sorted sequence. Let
$IQR=Q_3-Q_1$ and the upper threshold $U=Q_3+kIQR$. PERFIN flags the
high-expense side when

\begin{equation}
  z_i\geq2{,}5\qquad\hbox{or}\qquad x_i>Q_3+1{,}5IQR.
\end{equation}

If $s=0$ then $z_i=0$; if $IQR=0$ then the IQR branch raises no flag. The
output comprises the value, the date or label, $z_i$, the ratio to the mean
and the method \texttt{z}, \texttt{iqr} or \texttt{z+iqr}. The result only
asks the user to review; it does not prove fraud and the thresholds need to
be calibrated on representative data.

\subsection{Estimating Cash-Flow Runway}

Let $B=\sum_w B_w$ be the current total balance, $W=14$ the number of
observed calendar days and $e_j\geq0$ the spending on day $j$; days with no
spending have $e_j=0$. The burn rate and the number of remaining days are

\begin{equation}
  \bar e_W=\frac{1}{W}\sum_{j=1}^{W}e_j,\qquad
  d=\left\lfloor\frac{B}{\bar e_W}\right\rfloor.
\end{equation}

If $B\leq0$, \texttt{daysLeft=0}; if $\bar e_W=0$, the system does not
forecast a depletion date (\texttt{null}). Otherwise, the depletion date
equals the current date plus $d$. When there is a payday, the system finds
its next occurrence in the calendar and compares it with the depletion
date. For example, $B=2{,}8$ million VND and a 14-day total expense of
$1{,}4$ million VND give $\bar e=100{,}000$ VND/day, so $d=28$ days. This
is an extrapolation that keeps the recent spending level constant; it does
not yet account for upcoming income or an abnormal expense.

\subsection{Mining Recurring Expenses}

Only \texttt{expense} transactions not exceeding 500,000 VND are considered
in the current configuration. After the descriptions are normalized,
transactions with the same key are grouped into a cluster of $m$ amounts
$a_i$ and dates $t_i$. The average amount, relative deviation and average
cadence are

\begin{equation}
  \bar a=\frac{1}{m}\sum_{i=1}^{m}a_i,\qquad
  \delta_i=\frac{|a_i-\bar a|}{\bar a},\qquad
  \bar\Delta=\frac{1}{m-1}\sum_{i=2}^{m}|t_i-t_{i-1}|.
\end{equation}

A cluster is stable when every $\delta_i\leq0{,}15$. A cluster is a
recurring candidate when it is stable and satisfies one of two conditions:
$20\leq\bar\Delta\leq40$ days, or at least three occurrences. With $m=2$,
the cadence must still lie within the window. The monthly estimate is
currently $\bar a$ under the assumption of one occurrence per month; the
total subscription is the sum of the cluster estimates. Exact grouping
helps limit false positives but may miss variant descriptions. The result
is only a ``recurring signal''; it does not automatically create a
recurring bill or cancel a service.

\subsection{Cross-Category Spending Correlation}

For two weekly expense series $X=(x_1,\ldots,x_n)$ and
$Y=(y_1,\ldots,y_n)$ on the same weekly axis, a category's missing week is
filled with 0. The Pearson correlation coefficient~\cite{pearson1895} is

\begin{equation}
  r=\frac{\sum_i(x_i-\bar{x})(y_i-\bar{y})}
  {\sqrt{\sum_i(x_i-\bar{x})^2\sum_i(y_i-\bar{y})^2}}.
\end{equation}

The function returns 0 if $n<3$ or one series has zero variance. The
current version considers only categories with at least four weeks of
observation and reports only pairs with \textbf{positive} correlation
$r\geq0{,}6$; negative correlation is not yet reported. The system selects
the pair with the largest $r$ and must interpret it as ``co-varying in the
observed data'', not infer a causal relationship.

\subsection{Budget Recommendation and Forecasting}

Let $s_{c,j}$ be the expense of category $c$ in month $j$ over $m$ calendar
months; a month with no activity has $s_{c,j}=0$. The average and the base
level with a buffer $\beta=5\%$ are

\begin{equation}
  \bar s_c=\frac{1}{m}\sum_{j=1}^{m}s_{c,j},\qquad
  b_c=(1+\beta)\bar s_c.
\end{equation}

With monthly income $Y$, the default framework allots at most
$C_{need}=0{,}50Y$ to needs, $C_{want}=0{,}30Y$ to wants and sets a target
of $0{,}20Y$ for savings. Let $g$ be the spending group (needs or wants),
$C_g$ its allocation ceiling and $\sum_{k\in g}b_k$ the total base limit of
the categories in that group. Three strategies produce a raw limit $L_c$:

\begin{equation}
  L_c=
  \left\{
  \begin{array}{ll}
    b_c,
      & \mbox{category average},\\
    C_g\displaystyle\frac{b_c}{\sum_{k\in g}b_k},
      & \mbox{50/30/20},\\
    b_c\displaystyle\min\left(1,\frac{C_g}{\sum_{k\in g}b_k}\right),
      & \mbox{hybrid}.
  \end{array}
  \right.
\end{equation}

The result is rounded to a step of
10,000 VND; a positive value smaller than one step still becomes 10,000
VND. Confidence follows the number of months in which the category had
activity: high when at least 6, medium when 3--5, low when below 3; every
recommendation also carries a warning if the observed history is under
three months.

To forecast a limit overrun within the month, with $S$ the amount spent
after $d_e$ days and a month of $D$ days:

\begin{equation}
  q=\frac{S}{d_e},\qquad \widehat S_D=qD,\qquad
  d_{exceed}=\left\lceil\frac{L-S}{q}\right\rceil.
\end{equation}

The system attaches \texttt{likely\_to\_exceed} when $\widehat S_D>L$. The
formula assumes a constant spending rate, so it is only an early warning;
any recommendation or limit change always requires the user's confirmation.

\subsection{Goal and Debt-Repayment Planning}

For a savings or purchase goal, let $T>0$ be the goal amount, $C\geq0$ the
current amount, $R=\max(0,T-C)$ the remaining shortfall and $M\geq0$ the
monthly contribution. If $M>0$:

\begin{equation}
  n=\left\lceil\frac{R}{M}\right\rceil.
\end{equation}

If the deadline is $n_D>0$ months away, the required contribution is
$M_D=\lceil R/n_D\rceil$, the monthly shortfall is $G=\max(0,M_D-M)$ and
the shortfall at the deadline is $H=\max(0,R-Mn_D)$. When $M=0$, the
completion time is undefined; when $R=0$, the goal is complete
immediately. What-if simply re-runs the same function with $M'=M+E$, where
$E$ is the freed-up spending, and does not modify the original plan.

For debt repayment, $L$ is the current outstanding balance, $i_a$ the
annual interest rate as a percentage, $r=i_a/(100\times12)$ the monthly
rate and $n$ the number of months. The level payment that completes on time
is

\begin{equation}
  P=
  \left\{
  \begin{array}{ll}
    L/n, & r=0,\\[2mm]
    \displaystyle\frac{Lr}{1-(1+r)^{-n}}, & r>0.
  \end{array}
  \right.
\end{equation}

The simulation uses the recurrence
$L_{k+1}=\max(0,L_k(1+r)-P)$ and adds $L_kr$ to the total interest. If
$P\leq Lr$ in the first month, the balance does not decrease and the system
returns the warning \texttt{NEGATIVE\_AMORTIZATION}. If the debt is not
settled within the default simulation limit of 600 months, the plan is not
feasible within the horizon. The formula lets FR-10 clearly separate the
computation from the LLM's interpretation.

\section{Evaluation Methodology}

\subsection{Extraction and Classification Accuracy}

For a field to be extracted, $TP$ is a correctly predicted value, $FP$ is a
wrongly predicted or spurious value and $FN$ is a ground-truth value that
was missed:

\begin{equation}
  Precision=\frac{TP}{TP+FP},\qquad
  Recall=\frac{TP}{TP+FN},\qquad
  F1=\frac{2\,Precision\,Recall}{Precision+Recall}.
\end{equation}

Exact match is achieved only when all required fields of a transaction are
correct together. Macro-F1 averages the F1 of each category so that
high-sample classes do not mask low-sample ones. These metrics must be
computed on an evaluation set independent of the development prompts or
rules.

\subsection{OCR and STT Error and Business Impact}

Here $S$, $D$, $I$ are the number of substitutions, deletions and
insertions respectively; $N$ is the number of characters or words in the
ground truth. The subscript $c$ denotes character-level counts (CER) and
$w$ denotes word-level counts (WER):

\begin{equation}
  CER=\frac{S_c+D_c+I_c}{N_c},\qquad
  WER=\frac{S_w+D_w+I_w}{N_w}.
\end{equation}

CER/WER measure raw text; transaction-field accuracy is what tells whether
a media error corrupted the amount, date, merchant or line item. That is
why a provider smoke test must not be presented as accuracy.

\subsection{Algorithm Correctness, Grounding and Performance}

Deterministic algorithms are checked with hand-computed expected values,
the absolute error $AE=|\hat y-y|$ and invariants such as the unchanged
total of a wallet transfer. For narration, let $N_{all}$ be the number of
numeric expressions in the response and $N_{supported}$ the number that map
to facts or an equivalent formatting operation:

\begin{equation}
  Numeric\ Faithfulness=\frac{N_{supported}}{N_{all}}.
\end{equation}
\begin{equation}
  Unsupported\ Rate=1-Numeric\ Faithfulness.
\end{equation}

This is a \textbf{target evaluation method}; a checker covering every
number is not yet fully implemented. Latency uses percentiles: p50 is the
median, p95 is the value that 95\% of samples do not exceed; text/media,
cache hit/miss must be separated and the provider, hardware and number of
runs recorded.

\section{The Controlled Supporting Role of the LLM}

\subsection{Function Calling Instead of Delegating Business Authority}

The transformer uses attention to model token relationships and is the
foundation of modern LLMs \cite{vaswani2017}. PERFIN leverages the ability
to understand varied phrasing, but does not use the model to compute the
ledger or statistics. Tools are declared by schema using the provider's
function-calling mechanism \cite{geminifunctioncalling}; the LLM proposes
the tool name and arguments, while the application validates, creates the
preview, awaits confirmation and only then executes. Structured output only
guarantees the JSON shape, not that the category is correct, that the two
wallets are distinct or that the date is within the valid domain.

\subsection{Responsibility Boundaries and Grounding}

\begin{table}[htbp]
  \centering
  \caption{Responsibility boundary between the deterministic core and the LLM}
  \label{tab:llm-responsibility}
  \small
  \begin{tabularx}{\textwidth}{|p{3.8cm}|p{4.0cm}|X|}
    \hline
    \textbf{Task} & \textbf{Responsible component} &
    \textbf{Rationale} \\ \hline
    Total income--expense, balance, limits & SQL and business service &
    Accurate and repeatable \\ \hline
    Trend, anomaly, runway, correlation & Deterministic algorithm functions
    & Have formulas and independent tests \\ \hline
    Budget and goals & Budget/Goal Planner & Have constraints and simulation
    \\ \hline
    Understanding dates, currency units, multiple intents & LLM or local
    parser & A language/context problem \\ \hline
    Selecting the tool and filling parameters & LLM, then the service checks
    & Bridge from utterance to API \\ \hline
    Writing the interpretation/persona & LLM or template fallback & Must not
    change the facts \\ \hline
    Write, delete, transfer, re-tag & Service after confirmation & Has data
    impact and needs audit \\ \hline
  \end{tabularx}
\end{table}

Insight facts must include the value, unit, observation window, sample
count, method and a missing-data warning. The persona may change the tone
in the spirit of behavioral design \cite{thaler2009} but does not change
the facts. The system currently separates the facts from the narrator and
has a template fallback; a comprehensive numeric checker is still a target
design that needs evaluation. The LLM does not connect directly to
PostgreSQL and does not execute tools on its own.

\subsection{Conditions for the Use of the LLM to Be Meaningful}

The LLM is only valuable if it delivers a measurable improvement over the
form and the local parser: higher field F1 or exact match, fewer
clarification rounds or actions, while latency, cost and the
unsupported-number rate stay within acceptable thresholds. An ablation must
use the same data and the same task. If it does not create a benefit
sufficient to offset the risk, the direct flow or the local parser is
preferred. This approach keeps the thesis focused on data and algorithms,
with the LLM as a controlled communication layer.

\section{Infrastructure, Technology and Selection Rationale}

\begin{table}[htbp]
  \centering
  \caption{Technologies, selection criteria and alternatives}
  \label{tab:technology-rationale}
  \small
  \begin{tabularx}{\textwidth}{|p{3.2cm}|X|X|}
    \hline
    \textbf{Technology} & \textbf{Why it fits the problem} &
    \textbf{Boundary or alternative} \\ \hline
    React Native + Expo & One codebase for text, image, audio input and
    preview on mobile & Contains no formulas; separate native only needed
    when a device-specific requirement arises \\ \hline
    Node.js + Express & I/O-oriented API; shares JavaScript with the pure
    functions and tests; suits a modular monolith & No need for the
    operational cost of microservices \\ \hline
    PostgreSQL \cite{postgresql2026} & Foreign keys, \texttt{DECIMAL},
    transactions and aggregate queries suit a ledger & Files or NoSQL struggle to guarantee the same
    level of constraints and rollback \\ \hline
    Redis \cite{redisexpire} & TTL, cache-aside, pending state and queue & Not a source of
    truth; the memory fallback is used only in development \\ \hline
    BullMQ \cite{bullmq2026} & Scheduling, retry and job IDs for idempotent
    proactive tasks &
    A simple cron lacks the same queue/retry mechanism \\ \hline
    Gemini or local parser & The LLM extends language coverage; the parser
    guarantees a testable fallback & The provider sits behind an adapter and
    does not lock the business logic \\ \hline
    Local or cloud OCR/STT & Allows comparison of accuracy, latency, cost
    and privacy & Raw output always passes through the same
    validation/preview \\ \hline
  \end{tabularx}
\end{table}

Redis/BullMQ serve recurring reminders, runway scans, subscription scans,
month-end reports and export cleanup. Jobs must be small and idempotent; a
unique event key in the internal message prevents a duplicate effect on
retry. Scheduled auto-backup does not yet have a corresponding worker and
must not be regarded as operational.

\section{Related Research and Applications}

Research on financial literacy points to the role of planning and
long-term decision-making~\cite{lusardi2014}. Linear
regression~\cite{montgomery2012}, z-score, IQR~\cite{tukey1977} and
correlation~\cite{pearson1895} are explainable techniques, suited to the
thesis goal because their output can be checked independently.
OCR~\cite{smith2007}, STT~\cite{radford2022} and multimodal
LLMs~\cite{gemini2023} extend the input channels but do not on their own
guarantee the correctness of a financial record.

Commercial financial-management systems typically fall into three groups.
The first is form/dashboard applications such as Money
Lover~\cite{moneylover} or MISA MoneyKeeper~\cite{misamoneykeeper}: they
have well-structured data and rich visual reports, but entry still follows
multi-step forms and offers little support for everyday Vietnamese natural
language. The second is financial chatbots: natural to talk to, but the
numbers they return are often hard to trace back to a source and prone to
drift. The third is deep analytic models: strong on statistics but lacking
a confirmation loop for the everyday user. PERFIN does not compete on
commercial completeness. The chosen gap is to combine a relational ledger,
deterministic algorithms, a human in the loop and a bounded language
layer.

\begin{table}[htbp]
  \centering
  \caption{Gaps and PERFIN's approach to addressing them}
  \label{tab:research-gaps}
  \small
  \begin{tabularx}{\textwidth}{|X|X|X|}
    \hline
    \textbf{Gap} & \textbf{Risk} &
    \textbf{Approach} \\ \hline
    Natural entry but easy to record incorrectly & Dirty data spreads into
    reports & Typed draft, validation, preview, confirm \\ \hline
    Media has raw text but the fields are not certainly correct & Wrong
    amount, date, category & Measure CER/WER and field accuracy; confirm
    before writing \\ \hline
    Chatbot returns numbers of unclear origin & Hallucination and hard to
    audit & Query tools, JSON facts, target grounding checker \\ \hline
    Dashboard has numbers but lacks interpretation & Hard to notice patterns
    & Deterministic algorithms; the narrator only reads facts \\ \hline
    Non-adaptive classification & Repeats the same error & Correction
    retrieval, fuzzy threshold and margin \\ \hline
    Retry creates duplicate notifications & Poor experience and wrong data &
    Job ID, unique event key, idempotent handler \\ \hline
  \end{tabularx}
\end{table}
